\documentclass[12pt,a4paper]{article}
\usepackage{fontspec}
\setmainfont{cmunrm.otf}[BoldFont=cmunbx.otf,ItalicFont=cmunti.otf,BoldItalicFont=cmunbi.otf]
\setmonofont{cmuntt.otf}[Scale=0.9]
\usepackage[russian]{babel}
\usepackage[left=30mm,right=15mm,top=20mm,bottom=20mm]{geometry}
\usepackage{graphicx}
\usepackage{float}
\usepackage{pdflscape}
\newcommand{\University}{Санкт-Петербургский политехнический университет Петра Великого}
\newcommand{\Institute}{Институт компьютерных наук и кибербезопасности}
\newcommand{\School}{Высшая школа программной инженерии}
\newcommand{\Discipline}{Методики разработки современных веб-приложений}
\newcommand{\StudentGroup}{5140903/50401}
\newcommand{\StudentName}{К.И. Забродин}
\newcommand{\SupervisorRole}{доцент ВШ ПИ}
\newcommand{\SupervisorName}{А.В. Иванищев}
\newcommand{\ReportCity}{Санкт-Петербург}
\newcommand{\ReportYear}{2026}

\linespread{1.3}
\setlength{\parindent}{1.25cm}
\setlength{\parskip}{0pt}
\setlength{\emergencystretch}{3em}
\sloppy
\newcommand{\shot}[3]{
  \begin{figure}[H]
    \centering
    \includegraphics[width=\linewidth,height=0.45\textheight,keepaspectratio]{images/#1}
    \caption{#2}
  \end{figure}
}
\begin{document}
\begin{titlepage}
  \thispagestyle{empty}
  \begin{center}
    \University\\
    \Institute\\
    \School
  \end{center}
  \vspace{30mm}
  \begin{center}
    \textbf{ЛАБОРАТОРНАЯ РАБОТА №2}\\[5mm]
    \textbf{Разработка API}\\
    \textbf{с использованием фреймворка Sinatra}\\[4mm]
    по дисциплине «\Discipline»
  \end{center}
  \vspace{47mm}
  \begin{flushright}
    \begin{minipage}{0.58\textwidth}
      Выполнил: студент гр. \StudentGroup\\
      \StudentName\\[5mm]
      Руководитель: \SupervisorRole\\
      \SupervisorName
    \end{minipage}
  \end{flushright}
  \vfill
  \begin{center}
    \ReportCity\\
    \ReportYear
  \end{center}
\end{titlepage}
\setcounter{page}{2}


\section{Цель и задание работы}

Цель работы --- изучить маршрутизацию Sinatra и подготовить основу API для сервиса рекрутинга.

Необходимо установить зависимости, создать приложение в классическом и модульном стилях, проверить прямой запуск и запуск через Rack. Требуется реализовать именованные параметры, маски splat, регулярное выражение, перенаправление и необязательную часть маршрута, а также подключить пространство имён. В завершение нужно подготовить структуру приложения Sinatra для дальнейшего подключения PostgreSQL.

Работа выполнена в каталоге \texttt{sinatra\_api} репозитория \texttt{github.com/kiriklzz/}\allowbreak\texttt{ruby\_api\_labs}. Использованы Ruby 3.2.2, Sinatra 4.2.1, Rack и Puma 8.0.2. В методичке приведены более ранние версии Sinatra и серверы Thin/WEBrick; в работе те же способы запуска проверены с Puma.

\section{Установка и структура проекта}

В \texttt{Gemfile} указаны Sinatra, расширение \texttt{sinatra-contrib}, Rackup и Puma. Для следующих этапов API подключены драйвер \texttt{pg}, ORM Sequel, \texttt{sequel-seed}, Rake и Dotenv. Версии зависимостей сохранены в \texttt{Gemfile.lock}.

\begin{verbatim}
cd sinatra_api
bundle install
\end{verbatim}

Сохранены отдельные файлы: \texttt{classic.rb} для классического стиля и \texttt{my\_sinatra.rb} для модульного. Конфигурации Rack находятся в \texttt{classic.ru} и \texttt{modular.ru}. Итоговое приложение собирается через \texttt{config.ru}.

\begin{verbatim}
sinatra_api/
  config/          # настройки подключения
  helpers/         # сериализация и общие методы
  models/          # модели Sequel
  routes/          # маршруты по ресурсам
  lib/tasks/       # задачи Rake
  db/migrations/   # миграции
  seeds/           # исходные данные
  test/            # проверки маршрутов и API
\end{verbatim}

Маршруты сгруппированы по ресурсам в \texttt{routes}; отдельные контроллеры для этой структуры не потребовались. Миграции, наполнение базы и конечные точки с фильтрами относятся к следующим работам; здесь проверяется основа приложения и загрузка файлов.

\clearpage
\section{Классический и модульный стили}
\subsection{Классический вариант}

В \texttt{classic.rb} подключён \texttt{sinatra}. Маршруты определены на верхнем уровне. Корневой маршрут возвращает приветствие, \texttt{/hello/:name} использует имя из URL, а общий \texttt{/*} обрабатывает остальные пути.

\begin{verbatim}
require 'sinatra'
get('/') { 'Hello My Sinatra - Easy and Wide world!' }
get('/hello/:name') do
  "Sinatra приветствует тебя, #{params[:name]}!"
end
\end{verbatim}

Прямой запуск выполняется командой \texttt{bundle exec ruby classic.rb}. Для Rack создан \texttt{classic.ru}:
\begin{verbatim}
require_relative 'classic'
run Sinatra::Application
\end{verbatim}

Проверены оба варианта. На рисунке~1 приведён протокол прямого запуска и реальные ответы. Для параллельной проверки использован порт 4569; стандартный порт Sinatra --- 4567.

\shot{01_launch.png}{Прямой запуск классического приложения и HTTP-ответы}{65mm}

\clearpage
\subsection{Модульный вариант и запуск через Rack}

В \texttt{my\_sinatra.rb} создан класс \texttt{MyAppModule::App}, унаследованный от \texttt{Sinatra::Base}. Расширение пространства имён зарегистрировано в классе приложения.

\begin{verbatim}
require 'sinatra/base'
require 'sinatra/namespace'
module MyAppModule
  class App < Sinatra::Base
    register Sinatra::Namespace
    # определения маршрутов
    run! if app_file == $PROGRAM_NAME
  end
end
\end{verbatim}

Файл \texttt{modular.ru} подключает класс и передаёт его Rack:
\begin{verbatim}
require_relative 'my_sinatra'
run MyAppModule::App
\end{verbatim}

Выполнены запуск через Ruby-файл и запуск через Rack:
\begin{verbatim}
bundle exec ruby my_sinatra.rb
bundle exec rackup -p 4567 modular.ru
\end{verbatim}

Это альтернативные команды: предыдущий сервер нужно остановить либо выбрать другой порт. При контрольном запуске модульного Ruby-файла порт задан через \texttt{PORT=4571}. Для Rack порт задаётся аргументом \texttt{-p}.

На следующих рисунках показаны скриншоты протоколов фактических HTTP-запросов к модульному приложению. Протоколы содержат URL, статус, тип содержимого и тело ответа; это не снимки интерфейса Postman. Исходные результаты сохранены в \texttt{evidence/requests.json}.

\subsection{Порядок маршрутов}

Специальные маршруты объявлены до общего \texttt{/*}. Иначе общий обработчик перехватывал бы запросы к приветствию и другим примерам. Маска внутри namespace расположена перед внешней общей маской.

Текстовые примеры возвращаются с типом \texttt{text/plain; charset=utf-8}. Поэтому имя отображается как текст и не интерпретируется браузером как HTML.

\clearpage
\section{Проверка маршрутов}
\subsection{Именованные параметры}

Запрос \texttt{/hello/Peter} получает имя через \texttt{params[:name]}. Маршрут \texttt{/hello\_block/:name} передаёт параметр аргументом Ruby-блока. В \texttt{/players/:sport\_id} извлекается идентификатор спорта. Маршрут \texttt{/players/*/:year} сочетает splat с параметром года.

\shot{02_named.png}{Именованные параметры и сочетание с маской}{72mm}

\clearpage
\subsection{Namespace}

Внутри \texttt{namespace '/api/v1'} определён маршрут \texttt{/*}. У запроса \texttt{/api/v1/friends} параметр splat содержит только \texttt{friends}. Запрос \texttt{/friends} обрабатывается внешней маской. Корневой маршрут возвращает своё приветствие.

\shot{03_namespace.png}{Корневой маршрут и маски внутри и вне namespace}{62mm}

\clearpage
\subsection{Перенаправление}

Маршрут \texttt{/redirect} выполняет \texttt{redirect to('/hello/World')}. Автоматический тест отдельно проверяет наличие перенаправления, затем следует по адресу назначения. HTTP-клиент протокола также следует за перенаправлением; на рисунке показан конечный URL и итоговый ответ HTTP 200.

\shot{04_redirect.png}{Переход на маршрут приветствия после перенаправления}{50mm}

\subsection{Регулярное выражение и splat}

Регулярное выражение извлекает слово после \texttt{/meta/hello/}. Запрос \texttt{/meta/hello/12345/world} возвращает \texttt{Hello, 12345!}. В \texttt{/say/*/to/*} извлекаются две строки, а в \texttt{/download/*.*} --- путь и расширение файла. Массивы представлены JSON-строкой.

\shot{05_patterns.png}{Регулярное выражение, две маски и расширение файла}{72mm}

\clearpage
\subsection{Необязательный формат}

Шаблон \texttt{/jobs.?:format?} обслуживает запросы \texttt{/jobs}, \texttt{/jobs.json} и \texttt{/jobs.xml}. Во всех случаях возвращается одна демонстрационная строка. Суффиксы показывают работу шаблона и не переключают сериализацию ответа.

\shot{06_format.png}{Запросы с отсутствующим и заданным суффиксом формата}{65mm}

\clearpage
\section{Подготовка итогового API}

В \texttt{api\_sinatra\_postgres.rb} используется классический стиль. После подключения конфигурации базы загружаются файлы из \texttt{models}, \texttt{helpers} и \texttt{routes}. Пути строятся относительно \texttt{\_\_dir\_\_}, поэтому загрузка не зависит от рабочей папки терминала.

Файл \texttt{config.ru} запускает \texttt{Sinatra::Application}. Итоговое приложение проверено на порту 4568 одновременно с демонстрацией на 4567. Корневой запрос вернул \texttt{Hi! I am api\_sinatra\_postgres!} со статусом HTTP 200.

\shot{07_api.png}{Корневой ответ и файлы сборки итогового приложения}{65mm}

\clearpage
\section{Автоматическая проверка и результаты}

Для части 2 создан \texttt{test/routes\_test.rb}. Используются Minitest и Rack::Test. Проверены параметры, перенаправление, текстовый тип ответа, маски, регулярное выражение, необязательный формат и namespace.

\begin{verbatim}
bundle exec ruby -Itest test/routes_test.rb
\end{verbatim}

Результат отдельного запуска: 3 теста, 18 утверждений, 0 ошибок и 0 неудачных проверок. Прямые запуски и конфигурации Rack дополнительно проверены реальными HTTP-запросами. Общий набор частей 2--4 также прошёл: 10 тестов и 83 утверждения.

\shot{08_tests.png}{Результат тестов маршрутизации и версии среды}{80mm}

Настройка Run Configuration в RubyMine описана в README проекта. Запуск через эту IDE не выполнялся. Проверка проведена командами Ruby/Rack и HTTP-клиентом; снимки не выдаются за работу в RubyMine или Postman.

\section{Вывод}

Созданы классический и модульный варианты приложения Sinatra. Проверены прямой запуск и запуск через Rack, извлечение параметров URL, маски, регулярное выражение, перенаправление и namespace. Подготовлена структура итогового API с загрузкой моделей, общих методов и маршрутов. Результаты запросов и автоматические тесты подтверждают работу рассмотренных механизмов маршрутизации.

\end{document}
